\documentclass[a4paper]{article}     
\usepackage{amsfonts}
\usepackage{amsmath}
\usepackage{amssymb}
\usepackage{graphicx}
\usepackage{cancel}

\newcommand{\asa}{\Leftrightarrow} %als slechts als
\newcommand{\adR}{\Rightarrow} %als dan Rechts
\newcommand{\adL}{\Leftarrow}
\newcommand{\Lh}{\left(} %Links haakje
\newcommand{\Rh}{\right)}
\newcommand{\Lv}{\left[} %Links vierkanthaakje
\newcommand{\Rv}{\right]}
\newcommand{\La}{\left\{} %Links acolade
\newcommand{\Ra}{\right\}}
\newcommand{\Ras}{\right.} %Rechts acolade stelsel (omdat om stelsels te maken heb je een punt nodig
\newcommand{\pnR}{\rightarrow} %pijl naar Rechts
\newcommand{\limiet}{\lim\limits_}
\newcommand{\schrap}{\cancel}
\newcommand{\schrapb}{\bcancel} %backslash
\newcommand{\schrapk}{\xcancel} %kruis
\newcommand{\vervang}{\cancelto}
\newcommand{\intgrl}{\displaystyle\int} %integraal teken (bij het berekenen van primitieven)

\begin{document}

\noindent \large Auteur: Yoren Collier \\
\noindent \large Studierichting: Informatica\\
\noindent \large Oefeningenreeks 5A nr. 5.2\\

\medskip

\normalsize

$ BI \intgrl_{-2}^2 \Lh 4x^2 -5x +2 \Rh dx \adR OBI = \intgrl \Lh 4x^2 -5x +2 \Rh dx$\\

$= 4\intgrl x^2 dx -5 \intgrl x dx +2 \intgrl dx = \dfrac{4 x^3}{3} +c -\dfrac{5x^2}{2} +c +2x +c = \dfrac{4}{3} x^3 - \dfrac{5}{2} x^2 +2x +c $\\

$BI = \Lv \dfrac{4}{3} x^3 - \dfrac{5}{2} x^2 +2x +c \Rv_{-2}^2 = \dfrac{4}{3} \Lh 2 \Rh^3 - \dfrac{5}{2} \Lh 2 \Rh^2 + 2 \cdot 2 \schrap{+c} - \dfrac{4}{3} \Lh -2 \Rh^3 + \dfrac{5}{2} \Lh -2 \Rh^2 -2 \cdot \Lh -2 \Rh \schrap{-c} = \dfrac{32}{3} - \schrap{\dfrac{20}{2}} +4 + \dfrac{32}{3} + \schrap{\dfrac{20}{2}} +4 = \dfrac{64}{3} +8 = \dfrac{88}{3}$\\

\begin{thebibliography}{999}
%\bibitem{a} Referentie
\end{thebibliography}
\end{document} 